\documentclass[11pt]{article}

\usepackage{amsmath,amssymb,amsthm,mathtools}
\usepackage[margin=1in]{geometry}
\usepackage{microtype}
\usepackage[hidelinks]{hyperref}

\newtheorem{theorem}{Theorem}[section]
\newtheorem{lemma}[theorem]{Lemma}
\newtheorem{corollary}[theorem]{Corollary}
\newtheorem{proposition}[theorem]{Proposition}
\newtheorem{remark}[theorem]{Remark}

\newcommand{\dd}{\,\mathrm{d}}
\newcommand{\one}{\mathbf{1}}
\newcommand{\ip}[2]{\left\langle #1,#2\right\rangle}

\title{Rooted Triangle--Tree Extensions:\\
Two Positive Families for the Chromatic-Polynomial Graphon Bound}
\author{}
\date{}

\begin{document}

\maketitle

\begin{abstract}
For a finite graph $H$, put
\[
\Phi_H(p)=(1-p)^{v(H)}
\chi_H\!\left(\frac{1}{1-p}\right),
\]
with the value at $p=1$ understood by continuity.  We prove the bound
$t(H,W)\geq\Phi_H(p)$, throughout the full required interval
$p\geq 1/2$, for two infinite families.  The first is obtained from a
triangle by attaching any number of leaves to one triangle vertex.  The
second is obtained by attaching any number of pairwise disjoint two-edge
tails to one triangle vertex.  The proof is valid for every graphon and
uses a weighted rooted-triangle inequality and a separate one-variable
convexity argument.  In the six-vertex catalogue the new cases are Atlas
34, 36, and 92.
\end{abstract}


%%%%%%%%%%%%%%%%%%%%%%%%%%%%%%%%%%%%%%%%%%%%%%%%%%%%%%%%%%%%%%%%%%%%%%%%%%%%%%%
\section{Graphons and the two graph families}
%%%%%%%%%%%%%%%%%%%%%%%%%%%%%%%%%%%%%%%%%%%%%%%%%%%%%%%%%%%%%%%%%%%%%%%%%%%%%%%

Let $(\Omega,\mu)$ be a probability space.  A graphon is a symmetric
measurable function $W\colon\Omega^2\to[0,1]$.  For a finite simple graph
$H$, its homomorphism density is
\[
t(H,W)
=
\int_{\Omega^{V(H)}}
\prod_{uv\in E(H)}W(x_u,x_v)
\prod_{u\in V(H)}\dd\mu(x_u).
\]
No injectivity is imposed on the maps represented by this integral.  Write
\[
p=t(K_2,W)
\quad\text{and}\quad
d(x)=\int_\Omega W(x,y)\dd\mu(y).
\]
Thus $0\leq d\leq1$ and $\int d=p$.

Fix a triangle on vertices $o,a,b$, with $o$ distinguished as its root.

\begin{itemize}
\item For an integer $r\geq0$, let $L_r$ be obtained by adding vertices
$u_1,\ldots,u_r$ and precisely the additional edges $ou_i$.  Thus $L_r$
is a triangle with $r$ leaves having the same neighbour $o$.

\item For an integer $s\geq0$, let $Q_s$ be obtained by adding vertices
$y_1,z_1,\ldots,y_s,z_s$ and precisely the additional edges
$oy_i,y_iz_i$.  Thus $Q_s$ is a triangle with $s$ internally disjoint
two-edge tails having the same initial vertex $o$.
\end{itemize}

In particular, $L_0=Q_0=K_3$, while $L_1$ is the paw graph.

Every new vertex belongs to a tree attached to the triangle.  Once the
triangle is properly coloured, each new vertex has exactly $x-1$ choices
after the colour of its parent has been fixed.  Hence
\begin{align}
\chi_{L_r}(x)
&=x(x-1)^{r+1}(x-2),
&v(L_r)&=r+3,
\label{eq:chrom-L}
\\
\chi_{Q_s}(x)
&=x(x-1)^{2s+1}(x-2),
&v(Q_s)&=2s+3.
\label{eq:chrom-Q}
\end{align}
Both families have chromatic number $3$.  If $q=1-p$, direct substitution
in the definition of $\Phi$ gives
\begin{align}
\Phi_{L_r}(p)
&=
q^{r+3}\frac1q
\left(\frac pq\right)^{r+1}
\left(\frac{2p-1}{q}\right)
=p^{r+1}(2p-1),
\label{eq:target-L}
\\
\Phi_{Q_s}(p)
&=
q^{2s+3}\frac1q
\left(\frac pq\right)^{2s+1}
\left(\frac{2p-1}{q}\right)
=p^{2s+1}(2p-1).
\label{eq:target-Q}
\end{align}
These polynomial identities also hold at $p=1$ by continuity.


%%%%%%%%%%%%%%%%%%%%%%%%%%%%%%%%%%%%%%%%%%%%%%%%%%%%%%%%%%%%%%%%%%%%%%%%%%%%%%%
\section{A weighted rooted-triangle inequality}
%%%%%%%%%%%%%%%%%%%%%%%%%%%%%%%%%%%%%%%%%%%%%%%%%%%%%%%%%%%%%%%%%%%%%%%%%%%%%%%

Define the rooted triangle density
\[
\tau(x)
:=
\int_{\Omega^2}
W(x,y)W(x,z)W(y,z)
\dd\mu(y)\dd\mu(z).
\]

\begin{lemma}[Weighted rooted Goodman inequality]
\label{lem:weighted-rooted}
Let $W$ have edge density $p>0$.  For every integer $k\geq0$,
\[
\int_\Omega d(x)^k\tau(x)\dd\mu(x)
\geq
\frac{2p-1}{p}
\int_\Omega d(x)^{k+2}\dd\mu(x).
\]
\end{lemma}

\begin{proof}
For $j\geq0$, put
\[
M_j=\int_\Omega d(x)^j\dd\mu(x).
\]
Let $T_W$ be the integral operator with kernel $W$, and set
\begin{align*}
I_k&=\int_\Omega d(x)^k\tau(x)\dd\mu(x),\\
J_k&=\int_{\Omega^2}d(x)^kW(x,y)d(y)
\dd\mu(x)\dd\mu(y)
=\ip{d^k}{T_Wd}.
\end{align*}
For $A,B\in[0,1]$ one has $AB\geq A+B-1$.  Apply this to
$W(x,y)$ and $W(x,z)$, multiply by the nonnegative quantity
$d(x)^kW(y,z)$, and integrate.  Tonelli's theorem gives
\begin{equation}
I_k\geq 2J_k-pM_k.
\label{eq:I-J}
\end{equation}

Put $U=1-W$, let $T_U$ be its integral operator, and write
\[
u=T_U\one=1-d,
\qquad q=1-p.
\]
Since $T_Ud=T_U(\one-u)=u-T_Uu$, we have
\[
T_Wd
=p\one-T_Ud
=d-q\one+T_Uu.
\]
Consequently,
\begin{equation}
J_k=M_{k+1}-qM_k+\ip{d^k}{T_Uu}.
\label{eq:J-decomp}
\end{equation}

We next prove
\begin{equation}
\ip{d^k}{T_Uu}
\geq
\int_\Omega d(x)^ku(x)^2\dd\mu(x).
\label{eq:opposite-order}
\end{equation}
By the symmetry of $U$, twice the difference between the left- and
right-hand sides of \eqref{eq:opposite-order} equals
\[
\int_{\Omega^2}
U(x,y)
\bigl(u(y)-u(x)\bigr)
\bigl(d(x)^k-d(y)^k\bigr)
\dd\mu(x)\dd\mu(y).
\]
This integrand is nonnegative: $d=1-u$, so $u$ and
$d^k=(1-u)^k$ are oppositely ordered.  This proves
\eqref{eq:opposite-order}.  Since $u=1-d$, equations
\eqref{eq:J-decomp} and \eqref{eq:opposite-order} yield
\begin{align*}
J_k
&\geq
M_{k+1}-qM_k
+\int d^k(1-d)^2
\\
&=
M_{k+2}-M_{k+1}+pM_k.
\end{align*}
Using this in \eqref{eq:I-J},
\[
I_k\geq 2M_{k+2}-2M_{k+1}+pM_k.
\]
Finally,
\begin{align*}
&2M_{k+2}-2M_{k+1}+pM_k
-\frac{2p-1}{p}M_{k+2}
\\
&\qquad=
\frac1p
\left(M_{k+2}-2pM_{k+1}+p^2M_k\right)
\\
&\qquad=
\frac1p\int_\Omega d(x)^k(d(x)-p)^2\dd\mu(x)
\geq0.
\end{align*}
The claimed inequality follows.
\end{proof}


%%%%%%%%%%%%%%%%%%%%%%%%%%%%%%%%%%%%%%%%%%%%%%%%%%%%%%%%%%%%%%%%%%%%%%%%%%%%%%%
\section{A local inequality for two-edge tails}
%%%%%%%%%%%%%%%%%%%%%%%%%%%%%%%%%%%%%%%%%%%%%%%%%%%%%%%%%%%%%%%%%%%%%%%%%%%%%%%

Set
\[
A(x):=(T_Wd)(x)=\int_\Omega W(x,y)d(y)\dd\mu(y).
\]
The quantity $A(x)$ is the rooted homomorphism density of a two-edge tail
whose initial vertex is fixed at $x$.  Notice that $0\leq A\leq1$ and
\begin{equation}
\int_\Omega A(x)\dd\mu(x)
=\int_\Omega d(x)^2\dd\mu(x)
\geq p^2.
\label{eq:mean-A}
\end{equation}
The equality is symmetry and Fubini; the inequality is Jensen.

\begin{lemma}[Pointwise rooted triangle bound]
\label{lem:pointwise}
For almost every $x\in\Omega$,
\[
\tau(x)\geq \max\{2A(x)-p,0\}.
\]
\end{lemma}

\begin{proof}
The inequality
$W(x,y)W(x,z)\geq W(x,y)+W(x,z)-1$, multiplied by
$W(y,z)\geq0$ and integrated in $y,z$, gives
\[
\tau(x)
\geq
2\int_\Omega W(x,y)d(y)\dd\mu(y)-p
=2A(x)-p.
\]
Also $\tau(x)\geq0$ by definition.  Taking the larger of the two lower
bounds proves the assertion.
\end{proof}

\begin{lemma}[Tail convexity]
\label{lem:tail-convexity}
Fix $p>0$ and an integer $s\geq0$.  The function
\[
f_{p,s}(a)=a^s\max\{2a-p,0\},
\qquad 0\leq a\leq1,
\]
where $a^0=1$, is convex and nondecreasing.
\end{lemma}

\begin{proof}
For $s=0$ this is the maximum of the two affine functions $0$ and
$2a-p$.  Suppose $s\geq1$.  The function vanishes on $[0,p/2]$, while
on $[p/2,1]$ it is
\[
2a^{s+1}-pa^s.
\]
Its first derivative on the latter interval is
\[
a^{s-1}\bigl(2(s+1)a-ps\bigr)>0,
\]
and its second derivative there is
\[
s a^{s-2}\bigl(2(s+1)a-p(s-1)\bigr)\geq0.
\]
At $a=p/2$ the right derivative is
$p(p/2)^{s-1}>0$, whereas the left derivative is $0$.  Thus the
one-sided derivatives increase across the joining point, proving both
convexity and monotonicity on the whole interval.
\end{proof}


%%%%%%%%%%%%%%%%%%%%%%%%%%%%%%%%%%%%%%%%%%%%%%%%%%%%%%%%%%%%%%%%%%%%%%%%%%%%%%%
\section{The positive theorem}
%%%%%%%%%%%%%%%%%%%%%%%%%%%%%%%%%%%%%%%%%%%%%%%%%%%%%%%%%%%%%%%%%%%%%%%%%%%%%%%

\begin{theorem}[Rooted triangle--tree extensions]
\label{thm:main}
Let $W$ be an arbitrary graphon with edge density $p$.  For every
$p\in[1/2,1]$ the following statements hold.

\begin{enumerate}
\item For every integer $r\geq0$,
\[
t(L_r,W)
\geq p^{r+1}(2p-1)
=\Phi_{L_r}(p).
\]

\item For every integer $s\geq0$,
\[
t(Q_s,W)
\geq p^{2s+1}(2p-1)
=\Phi_{Q_s}(p).
\]
\end{enumerate}

Since $\chi(L_r)=\chi(Q_s)=3$, the interval $p\geq1/2$ is exactly the
full required interval
\[
p\geq1-\frac1{\chi(H)-1}.
\]
\end{theorem}

\begin{proof}
Integrating out the $r$ leaves of $L_r$, conditional on the image $x$ of
the root $o$, gives
\[
t(L_r,W)=\int_\Omega d(x)^r\tau(x)\dd\mu(x).
\]
Lemma~\ref{lem:weighted-rooted}, followed by Jensen's inequality, gives
for $p\geq1/2$
\begin{align*}
t(L_r,W)
&\geq
\frac{2p-1}{p}\int_\Omega d(x)^{r+2}\dd\mu(x)
\\
&\geq
\frac{2p-1}{p}p^{r+2}
=p^{r+1}(2p-1).
\end{align*}
Here multiplication by $(2p-1)/p$ preserves the Jensen inequality
because this factor is nonnegative on the stated interval.

For $Q_s$, integrating out its $s$ disjoint two-edge tails gives
\[
t(Q_s,W)=\int_\Omega \tau(x)A(x)^s\dd\mu(x).
\]
By Lemma~\ref{lem:pointwise}, Lemma~\ref{lem:tail-convexity}, Jensen's
inequality, \eqref{eq:mean-A}, and monotonicity of $f_{p,s}$,
\begin{align*}
t(Q_s,W)
&\geq
\int_\Omega f_{p,s}(A(x))\dd\mu(x)
\\
&\geq
f_{p,s}\!\left(\int_\Omega A(x)\dd\mu(x)\right)
\\
&\geq f_{p,s}(p^2).
\end{align*}
For $p\geq1/2$ one has $p^2\geq p/2$, and therefore
\[
f_{p,s}(p^2)
=p^{2s}(2p^2-p)
=p^{2s+1}(2p-1).
\]
The target identifications are exactly \eqref{eq:target-L} and
\eqref{eq:target-Q}.
\end{proof}


%%%%%%%%%%%%%%%%%%%%%%%%%%%%%%%%%%%%%%%%%%%%%%%%%%%%%%%%%%%%%%%%%%%%%%%%%%%%%%%
\section{Endpoints and equality examples}
%%%%%%%%%%%%%%%%%%%%%%%%%%%%%%%%%%%%%%%%%%%%%%%%%%%%%%%%%%%%%%%%%%%%%%%%%%%%%%%

At the threshold $p=1/2$, both targets in
\eqref{eq:target-L}--\eqref{eq:target-Q} vanish, and the conclusion also
follows directly from nonnegativity of homomorphism densities.  Thus no
division by a vanishing quantity occurs.  At $p=1$, the condition
$0\leq W\leq1$ and $\int W=1$ forces $W=1$ almost everywhere, so both the
homomorphism density and the target equal $1$.

There are equality examples throughout the interval at all multipartite
grid points.  For an integer $k\geq2$, let $T_k$ be the balanced complete
$k$-partite graphon.  It has edge density $p_k=1-1/k$, and for every finite
graph $H$,
\[
t(H,T_k)=\frac{\chi_H(k)}{k^{v(H)}}=\Phi_H(p_k).
\]
In particular, this applies to every $L_r$ and $Q_s$.  For $k=2$ it gives
the threshold equality $0=0$.


%%%%%%%%%%%%%%%%%%%%%%%%%%%%%%%%%%%%%%%%%%%%%%%%%%%%%%%%%%%%%%%%%%%%%%%%%%%%%%%
\section{Catalogue consequences}
%%%%%%%%%%%%%%%%%%%%%%%%%%%%%%%%%%%%%%%%%%%%%%%%%%%%%%%%%%%%%%%%%%%%%%%%%%%%%%%

\begin{corollary}[New six-vertex catalogue cases]
\label{cor:atlas}
The following previously open Atlas graphs satisfy
$t(H,W)\geq\Phi_H(p)$ for every graphon $W$ throughout the full required
density interval $p\geq1/2$:
\[
\begin{array}{c|c|c|c}
\text{Atlas ID}&\text{graph6}&\text{family member}&\chi_H(x)\\
\hline
34&\texttt{D@\char123}&L_2&x(x-1)^3(x-2)\\
36&\texttt{DJc}&Q_1&x(x-1)^3(x-2)\\
92&\texttt{E?Fw}&L_3&x(x-1)^4(x-2)
\end{array}
\]
Their targets are, respectively,
\[
p^3(2p-1),\qquad p^3(2p-1),\qquad p^4(2p-1).
\]
\end{corollary}

\begin{proof}
Atlas 34 is a triangle with two leaves adjacent to the same triangle
vertex, so it is $L_2$.  Atlas 36 is a triangle with one two-edge tail,
so it is $Q_1$.  Atlas 92 is a triangle with three leaves adjacent to the
same triangle vertex, so it is $L_3$.  The graph6 codes give independent
machine-readable identifiers.  The conclusions follow from
Theorem~\ref{thm:main}.
\end{proof}


%%%%%%%%%%%%%%%%%%%%%%%%%%%%%%%%%%%%%%%%%%%%%%%%%%%%%%%%%%%%%%%%%%%%%%%%%%%%%%%
\section{Inputs and scope of the argument}
%%%%%%%%%%%%%%%%%%%%%%%%%%%%%%%%%%%%%%%%%%%%%%%%%%%%%%%%%%%%%%%%%%%%%%%%%%%%%%%

The proof above is self-contained apart from Tonelli--Fubini and Jensen's
inequality.  In particular, it does not invoke the odd-cycle theorem, the
pure-chordal theorem, clique--odd-cycle cases, universal-vertex lifting, an
arbitrary join closure, or any counterexample calculation.  The weighted
rooted-triangle inequality also appears as an ingredient in the separate
clique--odd-cycle note, but Lemma~\ref{lem:weighted-rooted} proves it here
in full.

All arguments use the ordinary (non-injective) graphon homomorphism
integrals directly.  They require neither a finite-step approximation nor
regularity of $W$, and they never replace the average edge-density
hypothesis by a minimum-degree or pointwise degree assumption.

%%%%%%%%%%%%%%%%%%%%%%%%%%%%%%%%%%%%%%%%%%%%%%%%%%%%%%%%%%%%%%%%%%%%%%%%%%%%%%%
\appendix
\section{What the Lean formalization actually proves}
%%%%%%%%%%%%%%%%%%%%%%%%%%%%%%%%%%%%%%%%%%%%%%%%%%%%%%%%%%%%%%%%%%%%%%%%%%%%%%%

Corollary~\ref{cor:atlas} is formalized: Atlas $34$, $36$ and $92$ are
\texttt{verified}, as is $L_1$ itself, the paw, which is Atlas $15$ and which
several later notes consume.  The development is
\texttt{lean/Taeyoung/Methods/RootedTriangleTree/}.  One difference of
substance.

\subsection*{Lemma~\ref{lem:tail-convexity} is not formalized}

The piecewise function $a\mapsto a^s\max\{2a-p,0\}$ does not occur, and neither
family goes through it.

For the leaf family $L_s$ the density factors as $\int d^s\tau$, and
Lemma~\ref{lem:weighted-rooted} sends that to $\frac{2p-1}{p}\int d^{s+2}$
directly; all that is then needed is $M_j\geq p^j$, which is Jensen for the
convex $t\mapsto t^j$ against the degree function.  The maximum with $0$ never
enters, because the weighted inequality has already been applied.

For the tail family $Q_s$, whose density is $\int\tau\cdot A$ with $A=T_Wd$ and
which does not factor as $\int d^r\tau$, the piecewise Jensen argument is
replaced by algebra.  The pointwise bound $\tau\geq2A-p$ of
Lemma~\ref{lem:pointwise} and $A\geq0$ give
\[
 \int\tau A\ \geq\ 2\!\int\!A^2-p\!\int\!A
 \ =\ 2M_2^2-pM_2+\Bigl(\int A^2-M_2^2\Bigr)
 \ \geq\ 2M_2^2-pM_2 ,
\]
the last step being Cauchy--Schwarz, and $m\mapsto2m^2-pm$ is increasing beyond
$m=p^2$ when $p\geq1/2$, so the bound descends to $2p^4-p^3=p^3(2p-1)$.  No
convexity of a piecewise function, and no one-sided derivatives at a joining
point.

\end{document}
